\chapter{Frontend Utility Classes}
\label{ch:fe-utilities}

Kiss provides a suite of JavaScript utility classes for date/time handling, number formatting, DOM manipulation, and general-purpose operations.  This chapter covers each class with examples.

\section{DateUtils}
\label{sec:feu-dateutils}

\texttt{DateUtils}\index{DateUtils.js} handles dates as YYYYMMDD integers --- the same format used by the Kiss backend.  This avoids timezone ambiguity and makes date arithmetic simple.

\begin{lstlisting}[style=javascript]
let today = DateUtils.today();              // e.g., 20260404
let formatted = DateUtils.intToStr4(today); // "04/04/2026"
let short = DateUtils.intToStr2(today);     // "04/04/26"
let sql = DateUtils.intToSQL(today);        // "2026-04-04"
// "Sat April 4, 2026"
let long = DateUtils.longFormat(today);
\end{lstlisting}

\subsection{Conversion}

\begin{lstlisting}[style=javascript]
// String to integer
let dt = DateUtils.strToInt("04/04/2026");  // 20260404
let dt2 = DateUtils.SQLtoInt("2026-04-04"); // 20260404

// Universal converter (accepts Date, string, or number)
// today as YYYYMMDD
let dt3 = DateUtils.toInt(new Date());

// Integer to JavaScript Date
let jsDate = DateUtils.intToDate(20260404);
\end{lstlisting}

\subsection{Date Arithmetic}

\begin{lstlisting}[style=javascript]
let next = DateUtils.intAddDays(20260404, 7);    // 20260411
let prev = DateUtils.intAddDays(20260404, -30);  // 20260305
let diff = DateUtils.daysDifference(20260404, 20260101); // 93
// 12
let months = DateUtils.monthsDifference(20260404, 20250404);
\end{lstlisting}

\subsection{Date Components}

\begin{lstlisting}[style=javascript]
let y = DateUtils.year(20260404);       // 2026
let m = DateUtils.month(20260404);      // 4
let d = DateUtils.day(20260404);        // 4
let dow = DateUtils.dayOfWeek(20260404);     // 6 (Saturday)
let name = DateUtils.dayOfWeekName(20260404); // "Saturday"
let mon = DateUtils.monthName(20260404);      // "April"
\end{lstlisting}

\subsection{Validation}

\begin{lstlisting}[style=javascript]
DateUtils.isValid(20260404);       // true
DateUtils.isValid(20261332);       // false (invalid month/day)
DateUtils.isSQLDate("2026-04-04"); // true
\end{lstlisting}

\begin{notebox}
\texttt{DateUtils} is locale-aware.  \texttt{strToInt()} and \texttt{intToStr4()} automatically detect whether the locale uses MM/DD/YYYY or DD/MM/YYYY format.
\end{notebox}

\section{TimeUtils}
\label{sec:feu-timeutils}

\texttt{TimeUtils}\index{TimeUtils.js} handles times as HHMM integers (e.g., 1430 for 2:30 PM).

\begin{lstlisting}[style=javascript]
let now = TimeUtils.current();              // e.g., 1430
let ampm = TimeUtils.format(1430);          // "2:30 PM"
// "02:30 PM" (zero-filled hour)
let zf   = TimeUtils.format(1430, true);
\end{lstlisting}

\subsection{Conversion and Parsing}

\begin{lstlisting}[style=javascript]
// Parse various formats
TimeUtils.strToInt("2:30 PM");   // 1430
TimeUtils.strToInt("14:30");     // 1430
TimeUtils.strToInt("1430");      // 1430

// Universal converter
TimeUtils.toInt(new Date());     // current time as HHMM
\end{lstlisting}

\subsection{Time Arithmetic}

\begin{lstlisting}[style=javascript]
TimeUtils.addHours(1430, 2);      // 1630
TimeUtils.addMinutes(1430, 45);   // 1515
TimeUtils.diff(1630, 1430);       // 120 (minutes)
TimeUtils.totalMinutes(1430);     // 870
TimeUtils.minutesToTime(870);     // 1430
\end{lstlisting}

\subsection{Time Components}

\begin{lstlisting}[style=javascript]
TimeUtils.hours(1430);    // 14
TimeUtils.minutes(1430);  // 30
\end{lstlisting}

\section{DateTimeUtils}
\label{sec:feu-datetimeutils}

\texttt{DateTimeUtils}\index{DateTimeUtils.js} handles combined date/time values as YYYYMMDDHHMM integers and provides extensive timezone support.

\begin{lstlisting}[style=javascript]
// Current date/time
let now = DateTimeUtils.now();  // e.g., 202604041430

// Create from date + time
let dt = DateTimeUtils.create(20260404, 1430); // 202604041430

// Format for display (locale-aware)
DateTimeUtils.formatDate(202604041430);
    // "04/04/2026 2:30 PM"
\end{lstlisting}

\subsection{Epoch Millisecond Handling}

Backend datetime values often arrive as epoch milliseconds.  \texttt{DateTimeUtils} handles this automatically:

\begin{lstlisting}[style=javascript]
// Backend returns epoch millis (e.g., 1737292576000)
let formatted = DateTimeUtils.formatDate(res.createdAt);
    // "01/19/2025 11:36 AM"
\end{lstlisting}

\begin{warnbox}
Always use \texttt{DateTimeUtils.formatDate()} to format backend datetime values.  Never manually parse epoch milliseconds with substring operations --- that produces garbage like ``49/12/1768''.
\end{warnbox}

\subsection{Component Extraction}

\begin{lstlisting}[style=javascript]
let dt = 202604041430;
DateTimeUtils.getDate(dt);  // 20260404 (YYYYMMDD)
DateTimeUtils.getTime(dt);  // 1430 (HHMM)
\end{lstlisting}

\subsection{DateTime Arithmetic (DST-Aware)}

\begin{lstlisting}[style=javascript]
let dt = 202604041430;
DateTimeUtils.addDays(dt, 7);      // 202604111430
DateTimeUtils.addHours(dt, 3);     // 202604041730
DateTimeUtils.addMinutes(dt, 90);  // 202604041600
DateTimeUtils.addMonths(dt, 1);    // 202605041430
DateTimeUtils.addYears(dt, 1);     // 202704041430
\end{lstlisting}

\noindent
All arithmetic methods accept an optional timezone parameter and handle DST transitions correctly.

\subsection{Timezone Support}

\begin{lstlisting}[style=javascript]
// Get local timezone
DateTimeUtils.localZone();           // "America/New_York"
DateTimeUtils.getLocalTimezoneShortText(); // "EDT"

// Convert between timezones
DateTimeUtils.epochToDisplayEpoch(
    epochMs, 'America/New_York', 'America/Chicago');
\end{lstlisting}

\section{NumberUtils}
\label{sec:feu-numberutils}

\texttt{NumberUtils}\index{NumberUtils.js} provides number comparison and rounding:

\begin{lstlisting}[style=javascript]
// Approximate equality (useful for floating-point)
NumberUtils.floatEqual(0.1 + 0.2, 0.3, 0.0001); // true

// Rounding
NumberUtils.round(3.14159, 2);     // 3.14
NumberUtils.roundUp(3.141, 2);     // 3.15
NumberUtils.roundDown(3.149, 2);   // 3.14
\end{lstlisting}

\section{DOMUtils}
\label{sec:feu-domutils}

\texttt{DOMUtils}\index{DOMUtils.js} is the DOM abstraction layer.  It provides methods for element access, manipulation, visibility, attributes, events, and dimensions.

\begin{warnbox}
\textbf{Only code under the \texttt{kiss/} directory should use DOMUtils directly.}  Application code should use the Kiss component APIs (\texttt{\$\$()}, \texttt{getValue()}, etc.) instead.  This layered architecture ensures consistent DOM manipulation and easier maintenance.
\end{warnbox}

\noindent
Key DOMUtils categories:

\begin{description}[style=nextline]
  \item[Element access] \texttt{getElement(id)}, \texttt{query(selector)}, \texttt{queryAll(selector)}
  \item[Creation] \texttt{create(tag, attrs)}, \texttt{append(parent, html)}
  \item[Visibility] \texttt{show(el)}, \texttt{hide(el)}, \texttt{isVisible(el)}
  \item[Classes] \texttt{addClass(el, cls)}, \texttt{removeClass(el, cls)}, \texttt{toggleClass(el, cls)}
  \item[Attributes] \texttt{attr(el, name, val)}, \texttt{prop(el, name, val)}, \texttt{css(el, prop, val)}
  \item[Content] \texttt{text(el, text)}, \texttt{html(el, html)}, \texttt{val(el, value)}
  \item[Events] \texttt{on(el, event, handler)}, \texttt{off(el, event, handler)}
  \item[Navigation] \texttt{children(el)}, \texttt{parent(el)}, \texttt{closest(el, selector)}
  \item[Dimensions] \texttt{width(el)}, \texttt{height(el)}, \texttt{offset(el)}
  \item[Navigation protection] \texttt{preventNavigation(isActive, onBack)}
\end{description}

\subsection{Browser Back Button Prevention}

\texttt{DOMUtils.preventNavigation()}\index{DOMUtils.preventNavigation} prevents accidental back-button and page-reload navigation:

\begin{lstlisting}[style=javascript]
DOMUtils.preventNavigation(
    function () { return true; },  // active when returns true
    function () {                  // called on back button
        Utils.yesNo('Confirm',
            'Are you sure you want to logout?',
            function () { Server.logout(); });
    }
);
\end{lstlisting}

\begin{notebox}
Chrome blocks navigation prevention until the page has received a user gesture (click, keypress, or touch).  Before any interaction, the back button cannot be prevented.  After the first click or keypress, full protection activates.
\end{notebox}

\section{Utils}
\label{sec:feu-utils}

\texttt{Utils}\index{Utils.js} is the general-purpose utility class.  Its most important functions are covered in other chapters:

\begin{itemize}
  \item \texttt{Utils.loadPage()} --- Chapter~\ref{ch:frontend}
  \item \texttt{Utils.popup\_open()} / \texttt{popup\_close()} --- Chapter~\ref{ch:popups}
  \item \texttt{Utils.showMessage()} / \texttt{Utils.yesNo()} --- Chapter~\ref{ch:popups}
  \item \texttt{\$\$()} --- Chapter~\ref{ch:frontend}
\end{itemize}

\noindent
Additional utility methods:

\begin{longtable}{P{4.8cm}P{5.7cm}}
\toprule
\textbf{Method} & \textbf{Description} \\
\midrule
\endhead
\texttt{Utils.cleanup()} & Remove event handlers and state from the current screen before loading a new one. \\
\texttt{Utils.showReport(url)} & Open a report URL in a new browser tab. \\
\texttt{Utils.waitMessage(msg)} & Show a non-modal progress message. \\
\texttt{Utils.waitMessageEnd()} & Dismiss the progress message. \\
\texttt{Utils.rescan()} & Activate custom components in newly injected HTML. \\
\texttt{Utils.\allowbreak afterComponentsLoaded(\allowbreak fn)} & Execute a function after all components are loaded. \\
\texttt{Utils.\allowbreak useComponent(\allowbreak name)} & Register a UI component for use. \\
\texttt{Utils.getAppUrl()} & Get the application's base URL. \\
\texttt{Utils.saveData(key, val)} & Save a cross-screen data item by key.  Backed by \texttt{AppState}\index{AppState}, so per the configured store (Chapter~\ref{ch:routing-state}) it survives page navigation and reload; the value must be JSON-serializable. \\
\texttt{Utils.getData(key)} & Retrieve a cross-screen data item by key. \\
\texttt{Utils.\allowbreak getAndEraseData(\allowbreak key)} & Retrieve a cross-screen data item by key and remove it, returning its prior value. \\
\bottomrule
\end{longtable}

\section{MutableString}
\label{sec:feu-mutablestring}

\texttt{MutableString}\index{MutableString.js} provides efficient character-level string modifications:

\begin{lstlisting}[style=javascript]
let ms = new MutableString("Hello World");
ms.setCharAt(5, '-');          // "Hello-World"
ms.insert(5, " Beautiful");    // "Hello Beautiful-World"
ms.delete(5);                  // "Hell Beautiful-World"
ms.append("!");                // "Hell Beautiful-World!"
ms.replace("World", "Day");   // "Hell Beautiful-Day!"
let ch = ms.charAt(0);        // "H"
let len = ms.length();        // 20
let str = ms.toString();      // "Hell Beautiful-Day!"
ms.clear();                   // ""
\end{lstlisting}

\section{Editor}
\label{sec:feu-editor}

\texttt{Editor}\index{Editor.js} wraps CKEditor 5 for rich text editing:

\begin{lstlisting}[style=javascript]
// Create editor in a container div
const editor = await Editor.create('editor-container');

// Get/set content
let html = editor.getHtml();
let text = editor.getText();
editor.setHtml('<p>Hello <b>World</b></p>');
editor.clear();

// Read-only mode
editor.readOnly();
editor.readWrite();

// Print editor contents
editor.printEditorContent();

// Cleanup (usually handled by framework)
editor.dispose();
\end{lstlisting}

\noindent
Editors are automatically cleaned up when screens unload via the editor context management system, similar to grids.
